% s-226-p226-m002; original PDF 226, 227
\Aphorism{s-226-p226-m002}{318}
Als Nichtmaler möchte ich mir erlauben, in einer modernen Richtung der Malerei, die sich zur Fahne des \Sperr{plein} \Sperr{air} bekennt, einfach den Sieg der Undulationstheorie über die rein geometrische Optik zu erblicken. Man kann das Licht einerseits als blosses Mittel des Sichtbarwerdens der Körper nach Form und Farbe behandeln; so urtheilt der naive Realist, der die Dinge wirklich so beschaffen glaubt, wie er sie sieht — und ähnlich fasst die geometrische Optik nur dasjenige am Licht ins Auge, was geradlinige Überlieferung der leuchtenden, lichtreflectirenden oder lichtdurchlassenden Gebilde ist. Eine Malerei, die sich daran hält, wird auf correcte Zeichnung den Hauptwerth legen und die Farbe in scharf-getrennten Flächen geben; sie schätzt sozusagen das Licht am höchsten, von dem am wenigsten gesprochen wird, das am wenigsten als selbständige Realität ins Bewusstsein tritt. Dies gerade aber will umgekehrt die plein-air-Kunst; sie fasst den optischen Hergang nicht als reinen Bericht, sondern als freie Umschreibung auf, sie kennt das Licht als eigene Physis mit eigenen Gesetzen, die sich keineswegs zum blossen Objectiviren der geometrischen Umrisse und Linien herablässt, sondern selbstthätig ins Spiel mischt und ihre Phänomene zu Gesicht bringt. Das Licht quillt, schäumt über, trocknet ein, überall treten Beugungen, Farbensäume, Änderungen des Polarisationszustandes auf; unter den Oberflächenfarben hervor dringen Farben aus dem Inneren der Körper der Himmel ist nicht gleichmässig blau, sondern hat blaueste Stellen inmitten weniger blauer; die Sonnenstrahlen, die durch Laub fallen, zeichnen nicht einfach die Conturen der Blätter, sondern verwischen sie durch Flimmern und erregen an den Blatträndern neues Lichtspiel. Diese Art Malerei will das volle Licht, das von sich selbst und nicht nur von der Lichtquelle erzählt; sie sucht nicht die kürzeste und getreueste Relation zwischen Object und Auge, sondern eher die complicirteste, den an schönen Nebenerscheinungen reichsten Lichtweg — ihr ist das Licht selbst Object, während es jenen Anderen nur Mittel zum Zweck ist.
\par

